\documentclass[11pt]{article}
\usepackage[margin=3cm]{geometry}
\usepackage{graphicx}
	\usepackage{times}

\title{TurtleKV Recovery Design}
\author{
Tony Astolfi
\and
Gabe Bornstein
}

\begin{document}
\maketitle{}

%%------------------------------------------------------------------------------
\section{Concepts}\label{section:concepts}

\subsection{Edit}

A single-key modification to the database.

\subsection{EditOffset (bytes)}

A scalar quantity that uniquely identifies a location in the logical
ordering of edits to the database.  Equal to the total size of all the
individual edit records since the database was in an \textit{Empty}
state.

Every Edit is assigned a unique EditOffset interval
\texttt{$e_{begin}$...$e_{end}$}, when the edit is applied to the
database.  Edits to a given key must be applied in-order on recovery
(R7) for consistency, as the order of update may affect the eventual
key value.  We also need to recover the total order of update between
edits of different keys so that we can recover as many edits as
possible, and so that we can answer point-in-time queries consistently
across restart/recovery.

\subsection{Checkpoint}

A TurtleTree that captures the end state of all edits up to some
EditOffset, the \textit{UpperBound} of the Checkpoint.

\subsection{Log}\label{section:concepts:log}

An append-only, fixed maximum capacity file which stores edits more
recent than the latest committed Checkpoint.  The Log
``file\footnote{The Log may be a self-contained file, a contiguous
region of some other file, or a contiguous region of a raw block
device.}" begins with a 4KiB \textit{Meta-Block}, followed by a series
of fixed size (default: 8KiB) \textit{Blocks}:

\vspace{11pt}
\includegraphics[width=0.9\textwidth]{figures/LogLayout.pdf}
\vspace{11pt}

The Meta-Block contains a magic number, a version number (of the Log
file layout; \texttt{1} for the format described in this document),
the Block size $B$ and count $N$, and the approximate active upper
bound and lower bound Block index (the \textit{active range}).  The
active range marks the contiguous set of Blocks which are in active
use, within a configured number of bytes (default=32MiB; the active
range accuracy $\Delta$, given in units of Block size, is also stored
as a field in the Meta-Block).  During Normal Operation
(post-Recovery), the Log's Meta-Block is periodically updated to keep
the active range up to date, within the configured accuracy
($\Delta$).

Note, the ``upper-bound" of the active range may be less-than the
``lower-bound," since both are given as Block indices relative to the
Log file start; if this is the case, it indicates that the active
range of blocks wraps around from the end of the file (i.e. the last
Black) back round to Block 1.  The stored active range bounds are
approximate because the true bound may be anywhere from the stated
index $I$ up to $I+\Delta$.

The Log \textbf{must not} ever record an active range $I_{lower},
I_{upper}, \Delta$ such that (any of):

\iffalse{}
OK:
\begin{itemize}
	\item $I_{lower} \le I_{upper} < (I_{upper}+\Delta) \pmod{N} $
	\item $I_{upper} < (I_{upper}+\Delta) \le I_{lower} \pmod{N} $
	\item $(I_{upper}+\Delta) \le I_{lower} \le I_{upper} \pmod{N} $
\end{itemize}
\fi

\begin{itemize}
	\item[a.] $(I_{upper}+\Delta) < I_{upper} \le I_{lower} \pmod{N} $
	\item[b.] $I_{lower} < (I_{upper}+\Delta) \le I_{upper} \pmod{N} $
	\item[c.] $I_{upper} \le I_{lower} < (I_{upper}+\Delta) \pmod{N} $
\end{itemize}

\begin{figure*}
	\centering
	\includegraphics[width=0.98\textwidth]{figures/LogActiveRange.pdf}
	\caption{Valid and Invalid Active Log Ranges, as recorded by
          the Meta-Block.  $I_{lower}$ may not appear in between
          $I_{upper}$ and $I_{upper} + \Delta$ (with wrap-around)}
\end{figure*}

\subsection{MemTable}

An in-memory dictionary data structure that stores an index of Edits
newer than the most recent Checkpoint.  A given MemTable does not
necessary contain all Edits newer than a Checkpoint; there may be
several MemTables between Checkpoints.

The current MemTable implementation uses Adaptive Radix Trees (ARTs)
with Optimistic Lock Coupling (OLC).

\subsection{Block}

The Log is stored as a collection of fixed-size Blocks (currently of
size 8KiB).  Blocks are laid out (Figure \ref{figure:blocklayout}) as
a fixed-size header (64 bytes), followed by a series of
variable-length \textit{Slots}, and then finally a reverse-order array
of the ending offset of each Slot relative to the end of the header.
The Block header includes the union interval of the EditOffset
intervals of all edits contained in its Slots, stored as a
\texttt{u64}, the lower bound EditOffset (inclusive), and the 32-bit
delta of the upper bound (exclusive).  The delta between lower and
upper bound EditOffset in a Block must be less than $2^{32}$, because
we store the EditOffset of each slot as a \texttt{u32} offset from the
lower bound (\texttt{u64}) of the Block.

\begin{figure}
\centering
\includegraphics[width=0.25\columnwidth]{figures/BlockLayout.pdf}
\caption{Log Block layout}\label{figure:blocklayout}
\end{figure}

The order of Slots containing Edits in a Block \textbf{must} match the
order of their EditOffsets.\footnote{In practice, this comes for free
as a consequence of a design where Block buffers are allocated
exclusively to a single thread.}

%%------------------------------------------------------------------------------
\section{Requirements}\label{section:requirements}

The recovery design for TurtleKV is constrained by the following:

\begin{enumerate}
	\item[R1.] The Log must retain all edits newer than the most
          recent recovered Checkpoint (TurtleTree)
	\item[R2.] The Log must trim all edits that have been
          incorporated into a committed Checkpoint
	\item[R3.] The unit of writing to the Log is a Block
	\item[R4.] Blocks are fixed size (in the current design, 8KiB)
	\item[R5.] The Block size may be greater than the unit of
          atomic update for the Log's storage device (typically 512B
          up to 4KiB)
	\item[R6.] During Recovery, the system must be able to detect
          \textit{short writes} (i.e. Blocks which have only been
          partially written)
	\item[R7.] The original total order (EditOffset) of Edits must
          be preserved across recoveries
	\item[R8.] The Blocks in the Log must be written and
          trimmed/over-written in a consistent order, to minimize
          storage device overhead (i.e., the Block writing pattern
          should be ``circular")
	\item[R9.] There must be no Block in the Log whose EditOffset
          interval spans the upper bound offset of a Checkpoint (i.e.,
          all Blocks must be either after a Checkpoint, or included in
          it)
\end{enumerate}

When comparing two designs, all else being equal, the design of better
\textit{Quality}:

\begin{enumerate}
	\item[Q1.] Minimizes the stored footprint per item (space amplification)
	\item[Q2.] Minimizes memory overhead during normal (post-Recovery) operation
	\item[Q3.] Minimizes the latency of Recovery
	\item[Q4.] Minimizes the memory overhead of Recovery
\end{enumerate}

%%------------------------------------------------------------------------------
\section{Design}

\subsection{Normal Operation}

When an Edit is applied to the database, it is assigned a unique
EditOffset.  This is done by maintaining a per-database
\texttt{std::atomic<u64>} variable, \texttt{next\_offset},
representing the least unused EditOffset; when processing an edit $e$,
we compute the serialized Log Slot size $|e|$ of $e$, and then invoke:
\texttt{e\_offset = next\_offset.fetch\_add($|e|$)} to assign the
EditOffset of $e$.

Once the EditOffset of $e$ has been determined, $e$ is added to a Slot
in a Block, so it can be appended to the Log.  Each thread that uses
the database has its own exclusive Log Writer Context.  This Context
stores an atomic pointer to the current Block buffer for its owning
thread.  One (or more) background threads managed by the Log Writer
are responsible for periodically ``stealing" ownership of the Block
buffers so they can be written to the Log.  While ownership of a Block
buffer resides with the thread who created it, that thread has
exclusive access free from data races.

When a thread wishes to append an Edit to the Log, it requests a
destination buffer, which is a slice (corresponding to a Slot) of a
Block buffer from its Log Writer Context.  The inputs for this request
are:

\begin{itemize}
	\item \texttt{min\_lb}: The minimum required EditOffset lower
          bound for the Block
	\item \texttt{edit\_offset}: The EditOffset lower bound of the
          Edit
	\item \texttt{edit\_size}: The size $|e|$ of the serialized
          Edit (i.e., the size of the Slot)
\end{itemize}

The Block header and Slot sizes are updated to reflect the allocation
of a new Slot.  The Block's EditSlot upper bound is also increased (if
necessary) so that it includes the new slot:

\vspace{11pt} \small{\texttt{block.upper\_bound =
    max(block.upper\_bound, edit.lower\_bound + edit.size)}}
\vspace{11pt}

The Block that is returned to the caller (into which the edit is
serialized) will be guaranteed to have an EditOffset lower bound
not-before the specified minimum lower bound.  This lets the caller
guarantee that the Block into which a given Edit is serialized does
not also contain Edits older a given offset.

To guarantee a consistent order of updates between the MemTable index
and the stored EditOffsets associated with each slot, a writer only
requests a Block to write into once it has determined where the update
will go in the MemTable (ART) index.  At the point when an ART-node
write lock has been acquired by the writer thread, we do the
\texttt{fetch\_add} on \texttt{next\_offet}, bail out if the MemTable
has overflowed, and eventually obtain a writable slot in a Block
buffer.  Contention across threads should not be impacted by this
design since \texttt{fetch\_add} is non-blocking; because of this,
there is no possibility of deadlock introduced.

The Edit slot stores its EditOffset as a \texttt{u32} offset from the
lower bound (\texttt{u64}) of the Block.  In the unlikely event that
this would cause an integer overflow, we just retry the Block buffer
allocation after updating \texttt{min\_lb}.

\subsubsection{MemTable Full}

MemTables have a maximum capacity (bytes).  When an Edit is applied,
the assigned EditOffset may indicate that the Edit cannot be added to
the current MemTable without exceeding its capacity.  When this
happens, the MemTable is permanently locked against future
modifications (i.e., it is \textit{finalized}).

MemTable boundaries are candidates for Checkpoint (upper bounds), so
when a MemTable is finalized, the next MemTable must use a
\texttt{min\_lb} EditOffset value equal to \texttt{next\_offset} (when
the MemTable is created).

\subsubsection{Periodic Updating of \texttt{min\_lb}}\label{section:minlbrebase}

When the database is shut down and later recovered, the Checkpoint
distance parameter (which determines the capacity of a MemTable) may
be set to a smaller value than it was before the restart.  For
example, the memory budget for the database may have been decreased.
However, since (R9) requires that no Checkpoint may have an upper
bound that bisects the EditOffset interval of a Block, it may not be
possible to find a valid Checkpoint upper bound to fit the new
Checkpoint distance.

To prevent this from happening\footnote{or at least, to minimize the
negative impact of this effect}, MemTables should periodically
"re-base" their \texttt{min\_lb}, by increasing it to the current
\texttt{next\_offset}.  This may waste up to $B-c$ bytes in each
active Block buffer (where $B$ is the size of a Block and $c$ is some
small constant).  This overhead is amortized against the period of
re-base, which should not be less than the size of a Checkpoint update
batch ($=L$, the leaf size, by default 32MiB).

\iffalse{}
\subsubsection{Key \texttt{locator}s}

The current design provides two formats for Edit slots; one for when a
key is first seen by a MemTable, and one for when the key is updated
with a new value.  The former of these includes the key bytes in the
slot itself; the latter replaces the key with a 32-bit integer
\texttt{locator} field.  This is broken into the upper 16-bits, which
identifies the Block (index) within the scope of the MemTable which
created it, and the lower 16-bits, which is the slot index within that
Block.

In order to recover the association between a physical Block in the
Log and the MemTable-scoped Block index used in locators, the current
Block header contains a 64-bit value called \texttt{owner\_id}.  This
is problematic in the new recovery design, after the introduction of
Big MemTables, which decouple Checkpoint batches from MemTables in
order to eliminate the trade-off between higher in-memory costs for
queries on recent updates and lower update costs from larger
Checkpoint distance.

The problem is that Big MemTables change the association between
Blocks and MemTables from static and durable in storage to a dyanmic,
runtime relationship which is not saved and may change between
restarts.  Also, locators create a link between blocks which
constrains Log trimming: the Log cannot be trimmed at a given point if
doing so would discard a Slot referenced by a locator in some later
Block.

We could attempt to keep locators and introduce a notion of
\textit{Log Epochs}, which would be logical segments of the Log
delineated by the points at which the \texttt{min\_lb} is re-based as
described in $\S$\ref{section:minlbrebase}.  However, this makes the
job of the MemTable more complicated, as it now needs to track (in the
ART entries) whether a given key has been seen in the current epoch,
which blocks belong to which Epoch, etc.

It is not clear that the benefit of locators (reducing the size of
Edit Log Slots for frequently updated keys) outweigh these costs in
additional complexity and runtime overhead.  Therefore, this design
proposes to remove locators in favor of a single Edit slot layout,
which always includes the key bytes\footnote{There are several design
options available if we want to reintroduce something like this; e.g.,
we could only store keys the first time they appear within a Block.}.
\fi{}

%%------------------------------------------------------------------------------
\subsection{Recovery}

This document does not cover the design for recovering the latest
Checkpoint; that is already handled through our use of \texttt{LLFS}
Volumes to implement the Checkpoint Log (not to be confused with the
Write-Ahead-Log or Change-Log for recent updates, which is what the
unqualified term ``Log" refers to in this document).

\subsubsection{Inputs to Recovery}

\begin{itemize}
	\item The latest Checkpoint's EditOffset upper bound
	\item The current value of Checkpoint distance (in units of
          batch size = leaf page size)
	\item The set of recovered Block buffers, with their physical
          Log file offset (in units of Block size)
\end{itemize}

\subsubsection{Outputs/Effects of Recovery}

\begin{itemize}
	\item A single active MemTable
	\item An ordered sequence of zero or more finalized MemTables
	\item The Log file offsets (in units of Block size) of the
          oldest and newest active Blocks
	\item All MemTables contain ART entries for the keys edited
          within their time span (the MemTable's \textit{scope}), each
          with the latest value for that key within the MemTable's
          scope
	\item All MemTables contain owning references (i.e.,
          \texttt{boost::intrusive\_ptr} or equivalent) to all Blocks
          in their scope; ART entries will reference key/value data
          from these Blocks
	\item The \texttt{ChangeLogFile} object's (block) lower/upper
          bound members are initialized to \textit{logical}
          equivalents of the Log file offset bounds (see previous
          item): if the physical newest block offset is less than the
          oldest, its logical equivalent is physical offset plus the
          capacity of the Log; the logical lower bound is always the
          physical lower bound
	\item All Block buffers own a \texttt{ChangeLogReadLock}
          covering their logical offset in the Log
	\item The Checkpoint generator is initialized with the latest
          recovered Checkpoint upper bound
	\item The Checkpoint update pipeline is initialized with the
          finalized MemTables
\end{itemize}

\subsubsection{Recovery Algorithm}

\begin{figure*}
\centering
\includegraphics[width=0.85\columnwidth]{figures/RecoveryScenario.pdf}
\caption{\small{A possible recovery scenario.  Boxes represent Blocks;
    green checkmarks are recovered Blocks, red x's are unrecoverable
    Blocks; the left- and right-edges of each Block's box in the
    diagram represent the EditOffset interval of the Block. Cluster
    boundaries are shown by large-dashed gray vertical lines; there
    are three Clusters.  The blue vertical dotted line represents the
    upper bound of the longest contiguous sequence of recoverable Edit
    slots.  Everything to the right, in the gray hashed region, is not
    recoverable because of the top-most of the two Blocks with a red
    X.}}\label{figure:recoveryscenario}
\end{figure*}

\textbf{Phase 1: Recover Clusters from Blocks}

\begin{enumerate}
	\item Read all Blocks in the Log's active range
          ($\S$\ref{section:concepts:log}) as recovered set $R$
	\item Remove from $R$ any Blocks for which any of the following are true:
	\begin{enumerate}
		\item EditOffset upper bound is not greater-than the
                  latest recovered Checkpoint upper bound
		\item Magic number is not correct
		\item Data integrity hash does not match the actual
                  contents of the Block
	\end{enumerate}
	\item Partition $R$ into disjoint \textit{Clusters} ($C_1,
          C_2, C_3, ...$) such that any two Blocks $b_i, b_j$ are in
          the same Cluster if their EditOffset intervals overlap.
          Note, not every pair of Blocks in a given Cluster will
          necessarily have overlapping EditOffset intervals
		\begin{itemize}
		    \item[] In the example of Figure
                      \ref{figure:recoveryscenario}, Block $b_1$ must
                      be in the same cluster as $b_2$ because they
                      overlap, and the same is true for $b_2$ and
                      $b_3$; so, $b_1$ and $b_3$ end up in the same
                      Cluster $C_1$, even though their EditOffset
                      intervals are disjoint.
		\end{itemize}
	\item For all Clusters, verify that the total sizes of all
          Slots in all Blocks of that Cluster is exactly equal to the
          size of the EditOffset interval for the Cluster.  The first
          (in EditOffset order) Cluster for which this is \textit{not}
          true is the \textit{last recovered Cluster}; all Clusters
          after the last are dropped
	\item For the last recovered Cluster $C_{\Omega}$ (if there is
          one-- in the example, $C_3$):
	\begin{enumerate}
		\item Create a priority queue of all the recovered
                  Blocks in $C_{\Omega}$, prioritized by the
                  EditOffset of the first unprocessed Slot in the
                  Block
		\item While there are no gaps in the sequence of
                  recovered Slots, continue to pull the first Block
                  out of the queue, marking its first unprocessed slot
                  as processed (+recovered), and re-inserting the
                  Block into the queue
		\item When a gap is encountered (i.e., the lower bound
                  of the next unprocessed Slot is greater than the
                  upper bound of the previously processed Slot),
                  \textsc{Stop;} add the longest recovered prefix
                  $C_{\Omega}'$ of $C_{\Omega}$ to the set of
                  recovered Clusters
		\begin{itemize}
		    \item[] In the example of Figure
                      \ref{figure:recoveryscenario}, we stop when we
                      reach the offset of the first slot in $b_7$.  If
                      we continue past that point, to include the
                      remainders of $b_4, b_5,$ and $b_6$, there will
                      be some Edits missing (in unrecoverable block
                      $b_7$).  Therefore we only recover a prefix of
                      of the Slots/Edits in Blocks $b_4, b_5,$ and
                      $b_6$.  None of the Edits in $b_8$ are
                      recoverable, despite the fact that $b_8$ is
                      recoverable per se, because of the gap created
                      by $b_7$.
		\end{itemize}
	\end{enumerate}
\end{enumerate}

\textit{Phase 1 Outputs:}

\begin{itemize}
	\item Block Clusters: $C_1, C_2, ... C_{\Omega}'$
	\item Set the Log's true active range according to the
          intersection of the Meta-Block's approximate active range
          and the recovered Blocks:
	\begin{itemize}
		\item Recovered active lower bound = minimum index of
                  Blocks from the recovered Clusters that are in the
                  range $[I_{lower}, \min(I_{upper}, I_{lower} +
                    \Delta))$
		\item Recovered active upper bound = 1 + maximum index
                  of Blocks from the recovered Clusters that are in
                  the range $[I_{upper}, I_{upper}+\Delta)$
		\item (Note: this recovers the \textit{physical}
                  active range -- convert to logical active range by
                  adding $N$ (maximum configured Log Block count) to
                  the upper bound if it is less than the lower bound
	\end{itemize}
\end{itemize}

\textbf{Phase 2: Recover MemTables from Clusters}

\begin{enumerate}
	\item Create an empty MemTable
	\item For all clusters $\in \{C_1, C_2, ..., C_{\Omega}'\}$:
	\begin{enumerate}
	    \item If the next cluster can be added to the current
              MemTable without overflowing its maximum capacity
              (Checkpoint Distance), add all the slots in EditOffset
              order
	    \item Else finalize the current MemTable, create a new
              one, and try again
	\end{enumerate}
\end{enumerate}

\textit{Phase 1 Outputs:}

\begin{itemize}
\item Zero or more finalized MemTables, in increasing order of EditOffset range
\item A single active MemTable
\end{itemize}

\textbf{Phase 3: Recover Checkpoint Pipeline from MemTables}

\vspace{11pt} At this point the only thing that must be done is to
start up the Checkpoint Generator and Batch scanner threads, feeding
the finalized MemTables in order.  The recovered active MemTable
becomes the active MemTable for the KVStore.
\vspace{11pt}

%%------------------------------------------------------------------------------
\section{Alternate/Rejected Designs}

\subsection{Partial Ordering of Slots (with \texttt{u16} revision per slot)}

\subsubsection{Summary}

To save space in the log, only store 16bits of ``revision" number to
order updates of the same key.  This is based on the following
assumptions:

\begin{enumerate}
\item Slots within a block are ordered implicitly
\item Resolving the order of two updates with different keys is not
  important; what matters is same-key
\item Therefore we only need to store enough version information to
  order same-key updates in ``concurrent" Blocks, i.e. those whose
  EditOffset intervals overlap
\item 16bits is sufficient because we need 1-bit to resolve
  wrap-around (like TCP packet sequence numbers within a stream),
  there is at most 1 Block per thread at a time, and $2^{32}=32768$ is
  a safe upper bound on the number of concurrent writer threads
\end{enumerate}

\subsubsection{Reason for Rejecting}

Assumption (2.) above turns out not to be true for two reasons:

\begin{itemize}
\item[i.] We need a way to be able to determine the longest recovered
  contiguous sequence of edits in the \textit{last} Block cluster --
  otherwise we essentially are forced to roll-back everything since
  the last Checkpoint, which defeats the point of the Log
\item[ii.] Point-in-time queries (i.e., snapshot isolated reads)
  require being able to know exactly what was visible at a given
  EditOffset, down to the individual key update
\end{itemize}

\end{document}
